\section{Glossari de termes}

\begin{itemize}
\item \hypertarget{def:blockchain}{\textbf{\textit{Blockchain}}} Tecnologia de registre distribuït que permet emmagatzemar informació de manera immutable, transparent i verificable.

\item \hypertarget{def:smartcontract}{\textbf{\textit{Smart contract}}} Programa executat en una blockchain que automatitza processos i garanteix l’execució segura de regles predefinides.

\item \hypertarget{def:descentralitzacio}{\textbf{Descentralització}} Absència d’una autoritat central de control; les dades es distribueixen entre múltiples nodes.

\item \hypertarget{def:hash}{\textbf{\textit{Hash} criptogràfic}} Funció matemàtica que transforma dades en una empremta digital única i irreversible.

\item \hypertarget{def:claus}{\textbf{Clau pública / clau privada}} Sistema criptogràfic on la clau privada signa i la clau pública verifica.

\item \hypertarget{def:signaturadigital}{\textbf{Signatura digital}} Mecanisme criptogràfic que garanteix autenticitat i integritat d’un missatge.

\item \hypertarget{def:criptografiaasimetrica}{\textbf{Criptografia asimètrica}} Sistema criptogràfic basat en l’ús de dues claus: pública i privada.

\item \hypertarget{def:wallet}{\textbf{\textit{Wallet blockchain}}} Aplicació que permet gestionar claus criptogràfiques i interactuar amb la blockchain.

\item \hypertarget{def:adreca}{\textbf{Adreça \textit{blockchain}}} Identificador únic utilitzat per enviar o rebre transaccions.

\item \hypertarget{def:token}{\textbf{\textit{Token}}} Credencial digital utilitzada per autenticar o autoritzar operacions.

\item \hypertarget{def:transaccio}{\textbf{Transacció}} Operació registrada en blockchain que representa una acció, com ara l’emissió d’un vot.

\item \hypertarget{def:ledger}{\textbf{\textit{Ledger} distribuït}} Base de dades compartida i sincronitzada entre múltiples nodes sense control central.

\item \hypertarget{def:node}{\textbf{Node}} Equip connectat a la xarxa blockchain que participa en la validació i manteniment del registre distribuït.

\item \hypertarget{def:consensus}{\textbf{\textit{Consensus} (mecanisme de consens)}} Procés mitjançant el qual els nodes acorden l’estat vàlid de la blockchain.

\item \hypertarget{def:pow}{\textbf{\textit{Proof of Work (PoW)}}} Mecanisme de consens basat en la resolució de problemes computacionals complexos.

\item \hypertarget{def:pos}{\textbf{\textit{Proof of Stake (PoS)}}} Mecanisme de consens on la validació depèn de la participació econòmica dels nodes.

\item \hypertarget{def:gas}{\textbf{\textit{Gas} (Ethereum)}} Unitat que mesura el cost computacional necessari per executar operacions en la blockchain.

\item \hypertarget{def:immutable}{\textbf{Immutable}} Propietat que garanteix que una vegada registrades, les dades no poden ser modificades.

\item \hypertarget{def:zkp}{\textbf{\textit{Zero-Knowledge Proof (ZKP)}}} Protocol criptogràfic que permet demostrar una informació sense revelar-la.

\item \hypertarget{def:blindsignatures}{\textbf{Blind signatures}} Tècnica criptogràfica que permet signar informació sense revelar el seu contingut.

\item \hypertarget{def:ssi}{\textbf{Identitat autosobirana (SSI)}} Model d’identitat digital on l’usuari controla les seves credencials sense dependència d’una autoritat central.

\item \hypertarget{def:testnet}{\textbf{\textit{Testnet}}} Xarxa blockchain de proves utilitzada per desenvolupament sense valor econòmic real.

\item \hypertarget{def:ganache}{\textbf{Ganache}} Blockchain local utilitzada per desenvolupar i provar contractes intel·ligents.

\item \hypertarget{def:hardhat}{\textbf{Hardhat}} Entorn de desenvolupament per crear, provar i desplegar contractes intel·ligents.

\item \hypertarget{def:solidity}{\textbf{Solidity}} Llenguatge de programació utilitzat per escriure contractes intel·ligents en Ethereum.

\item \hypertarget{def:openzeppelin}{\textbf{OpenZeppelin}} Llibreria estàndard de contractes intel·ligents segurs i reutilitzables.

\item \hypertarget{def:vulnerability}{\textbf{\textit{Smart contract vulnerability}}} Error o debilitat en el codi d’un contracte intel·ligent que pot ser explotat.

\item \hypertarget{def:web3}{\textbf{Web3}} Conjunt de tecnologies que permeten la interacció amb aplicacions descentralitzades.

\item \hypertarget{def:dapp}{\textbf{\textit{DApp (Decentralized Application)}}} Aplicació que funciona sobre una xarxa blockchain sense servidor central.

\item \hypertarget{def:evoting}{\textbf{\textit{E-voting}}} Sistema electrònic que permet l’emissió i recompte de vots de manera digital.

\item \hypertarget{def:anonimat}{\textbf{Anonimat del vot}} Garantia que no es pot associar un vot amb la identitat del votant.

\item \hypertarget{def:verificabilitat}{\textbf{Verificabilitat}} Capacitat de comprovar que el vot ha estat registrat correctament sense comprometre l’anonimat.

\item \hypertarget{def:auditabilitat}{\textbf{Auditabilitat}} Capacitat de verificar de manera independent la integritat i correcció del sistema.

\item \hypertarget{def:frontend}{\textbf{\textit{Front-end}}} Part visible del sistema amb la qual interactua l’usuari.

\item \hypertarget{def:backend}{\textbf{\textit{Back-end}}} Component del sistema responsable de la lògica, autenticació i gestió de dades.

\item \hypertarget{def:apirest}{\textbf{API REST}} Interfície que permet la comunicació entre sistemes mitjançant peticions HTTP.

\item \hypertarget{def:jwt}{\textbf{\textit{JSON Web Token (JWT)}}} Estàndard per transmetre informació d’autenticació de forma segura.

\item \hypertarget{def:webauthn}{\textbf{\textit{WebAuthn}}} Estàndard d’autenticació segura basat en criptografia i dispositius biomètrics.

\end{itemize}